\chapter{مدیریت بسته‌های نرم‌افزاری}

چنانچه زمانی را در انجمن‌های کاربری لینوکس سپری کنید، بی‌شک با دیدگاه‌های متعددی درباره‌ی اینکه کدام توزیع لینوکس «بهترین» است، مواجه خواهید شد. اغلب این مباحث به موضوعات سطحی نظیر زیبایی‌شناسی محیط دسکتاپ یا رنگ‌بندی پیش‌فرض (مانند جبهه‌گیری‌های سلیقه‌ای علیه اوبونتو) محدود می‌شوند.

با این حال، از نگاه فنی، حیاتی‌ترین شاخص برای سنجش کیفیت و کارآمدی یک توزیع، سیستم مدیریت بسته‌ها و پویایی جامعه‌ی پشتیبان آن است. اکوسیستم نرم‌افزاری لینوکس بسیار سیال و متغیر است؛ به‌طوری که توزیع‌های پیشرو به‌طور معمول هر شش ماه یک نسخه‌ی جدید عرضه می‌کنند و بسته‌های نرم‌افزاری نیز به‌صورت روزانه به‌روزرسانی می‌شوند. برای مدیریت صحیح این حجم انبوه از تغییرات، وجود ابزارهای قدرتمند مدیریت بسته الزامی است.

مدیریت بسته‌ها به فرآیند نصب، به‌روزرسانی و نگهداری نرم‌افزارها بر روی سیستم اطلاق می‌شود. امروزه اکثر کاربران لینوکس نیازمندی‌های خود را مستقیماً از طریق مخازن رسمی (\en{Repositories}) توزیع مورد نظر تأمین می‌کنند. این رویکرد با دوران نخستین لینوکس که در آن کاربران ناچار بودند کد منبع (\en{Source Code}) را شخصاً دانلود و کامپایل کنند، تفاوت بنیادینی دارد. اگرچه امکان کامپایل دستی کد منبع همچنان به عنوان یکی از مزایای آزادی‌بخش دنیای متن‌باز پابرجاست، اما بهره‌گیری از بسته‌های از پیش کامپایل‌شده (\en{Pre-compiled Packages}) روشی به‌مراتب سریع‌تر و بهینه‌تر است.

در این فصل، به بررسی ابزارهای خط فرمان (\en{CLI}) جهت مدیریت بسته‌ها می‌پردازیم. با وجود اینکه توزیع‌های مدرن از رابط‌های گرافیکی (\en{GUI}) پیشرفته‌ای برای این منظور بهره می‌برند، تسلط بر ابزارهای خط فرمان ضروری است؛ چرا که این ابزارها قابلیت‌هایی را ارائه می‌دهند که اجرای آن‌ها در محیط‌های گرافیکی دشوار یا گاهی ناممکن است.

\section{اکوسیستم‌های مدیریت بسته}

توزیع‌های مختلف لینوکس از سازوکارهای متفاوتی برای مدیریت بسته‌های نرم‌افزاری بهره می‌برند و به طور کلی، یک بسته‌ی کامپایل شده برای یک توزیع خاص، با سایر توزیع‌ها سازگاری ندارد. بخش عمده‌ای از توزیع‌های لینوکس در یکی از دو دسته‌بندی بنیادین فناوری بسته‌بندی قرار می‌گیرند: خانواده دبیان (\en{Debian}) با پسوند فایل \file{.deb} و خانواده ردهت (\en{Red Hat}) با پسوند فایل \file{.rpm}. هرچند استثناهای قابل توجهی نظیر جن‌تو (\en{Gentoo})، اسلکور (\en{Slackware}) و آرچ لینوکس (\en{Arch Linux}) وجود دارند، اما اکثر توزیع‌های مشهور بر پایه‌ی یکی از این دو سیستم عمل می‌کنند که جزئیات آن‌ها در جدول زیر ارائه شده است.

\begin{table}[htbp]
\centering
\caption{خانواده‌های اصلی سیستم‌های مدیریت بسته}
\label{tab:pkg-families}
\begin{tabularx}{\textwidth}{l X}
\toprule
\textbf{سیستم مدیریت بسته} & \textbf{توزیع‌های شاخص (فهرست گزیده)} \\
\midrule
سبک دبیان (\file{.deb}) & دبیان، اوبونتو، لینوکس مینت، سیستم‌عامل رزبری پای (\en{Raspberry Pi OS}) \\
\addlinespace
سبک رد هت (\file{.rpm}) & فدورا، سنت‌اواس (\en{CentOS})، ردهت اینترپرایز لینوکس (\en{RHEL})، اوپن‌سوزه (\en{OpenSUSE}) \\
\bottomrule
\end{tabularx}
\end{table}

\section{سازوکار سیستم‌های مدیریت بسته}

در دنیای نرم‌افزارهای انحصاری (\en{Proprietary Software})، روش معمول توزیع نرم‌افزار این‌گونه است: یا باید یک رسانه‌ی فیزیکی مانند دیسک نصب تهیه کنید، یا به وب‌سایت شرکت سازنده مراجعه کرده و فایل نصبی را دانلود کنید؛ سپس با اجرای یک برنامه‌ی «نصب‌کننده» (\en{Installer}) آن را روی سیستم خود نصب کنید.

لینوکس رویکردی کاملاً متفاوت دارد. تقریباً تمام نرم‌افزارهایی که برای یک سیستم لینوکسی نیاز دارید، از طریق اینترنت در دسترس‌اند. بخش بزرگی از این نرم‌افزارها توسط توسعه‌دهندگانِ توزیع (\en{Distribution}) به‌صورت «بسته‌های نرم‌افزاری» (\en{Packages}) از پیش آماده و کامپایل‌شده عرضه می‌شوند که نصب آن‌ها بسیار ساده است. در کنار این‌ها، برخی نرم‌افزارها تنها به‌صورت کد منبع (\en{Source Code}) در دسترس هستند که برای استفاده از آن‌ها باید فرآیند نصب دستی، شامل کامپایل کردن کد، طی شود. جزئیات این فرآیند در فصل ۲۳ بررسی خواهد شد.

\subsection{کالبدشکافی فایل‌های بسته (Package Files)}

واحد بنیادین نرم‌افزار در سیستم‌های مدیریت بسته، «فایل بسته» است که در واقع یک آرشیو فشرده از تمامی مؤلفه‌های لازم برای نصب یک برنامه محسوب می‌شود. محتوای یک بسته به طور کلی به سه بخش اصلی قابل تقسیم است: نخست، فایل‌های اجرایی برنامه و داده‌های جانبی مورد نیاز آن که شامل فایل‌های باینری، کتابخانه‌های وابسته، فایل‌های پیکربندی و مستندات می‌شود. دوم، فراداده (\en{Metadata}) که اطلاعات توصیفی و مدیریتی بسته از جمله نام، نسخه، شرح عملکرد، لیست محتویات، و مهم‌تر از همه، پیش‌نیازهای نرم‌افزاری (وابستگی‌ها) را در بر می‌گیرد و نقش کلیدی در مدیریت وابستگی‌ها ایفا می‌کند. سوم، اسکریپت‌های پیکربندی که بسیاری از بسته‌ها برای انجام وظایف سیستمی پیش از استقرار فایل‌ها (\en{Pre-install}) یا پس از آن (\en{Post-install}) از آن‌ها بهره می‌برند.

فایل‌های بسته توسط شخصی به نام «نگهدارنده بسته» (\en{Package Maintainer}) تولید می‌شوند که معمولاً از اعضای تیم توسعه‌ی آن توزیع است. نگهدارنده، کد منبع را از توسعه‌دهنده اصلی (\en{Upstream}) دریافت کرده، آن را برای توزیع مورد نظر کامپایل می‌کند و فراداده‌ها و اسکریپت‌های نصب مربوطه را به آن می‌افزاید. در بسیاری از موارد، نگهدارنده وصله‌هایی (\en{Patches}) را روی کد منبع اصلی اعمال می‌کند تا از سازگاری کامل و یکپارچگی برنامه با سایر بخش‌های آن توزیع لینوکس اطمینان حاصل شود.

\subsection{مخازن نرم‌افزاری (Repositories)}

برخی پروژه‌های نرم‌افزاری ترجیح می‌دهند فرآیند بسته‌بندی (\en{Packaging}) و توزیع نرم‌افزار خود را به‌طور مستقل مدیریت کنند؛ اما در عمل، بخش عمده‌ی بسته‌های نرم‌افزاری توسط تیم‌های توسعه‌ی توزیع‌ها (\en{Distribution Maintainers}) یا نگهدارندگان شخص ثالث (\en{Third-party Maintainers}) تهیه و نگهداری می‌شوند. این بسته‌ها از طریق مخازن نرم‌افزاری (\en{Software Repositories}) در اختیار کاربران قرار می‌گیرند؛ یعنی پایگاه‌های داده‌ی متمرکزی که هزاران بسته را میزبانی می‌کنند و به‌طور اختصاصی برای یک توزیع خاص پیکربندی، بهینه‌سازی و نگهداری می‌شوند.

هر توزیع معمولاً چندین مخزن متناظر با مراحل مختلف چرخه‌ی توسعه‌ی نرم‌افزار ارائه می‌دهد. برای نمونه:

\begin{itemize}
  \item \textbf{مخزن پایدار (\en{Stable}):} حاوی بسته‌های آزمایش‌شده و آماده‌ی استفاده‌ی عمومی است.
  \item \textbf{مخزن آزمایشی (\en{Testing}):} برای کاربرانی در نظر گرفته شده که مایل‌اند نسخه‌های آتی را پیش از انتشار نهایی بررسی و باگ‌های احتمالی را گزارش کنند.
  \item \textbf{مخزن ناپایدار یا توسعه (\en{Unstable/Development}):} میزبان بسته‌هایی است که هنوز در حال توسعه‌اند و قرار است در نسخه‌ی پایدار بعدی توزیع گنجانده شوند.
\end{itemize}

علاوه بر مخازن رسمی، مخازن شخص ثالث (\en{Third-party Repositories}) نیز وجود دارند که خارج از ساختار رسمی توزیع مدیریت می‌شوند. این مخازن معمولاً میزبان نرم‌افزارهایی هستند که به دلایل حقوقی—از جمله محدودیت‌های مرتبط با حق ثبت اختراع (\en{Patent}) یا قوانین مدیریت حقوق دیجیتال (\en{DRM})—امکان گنجاندن مستقیم آن‌ها در مخازن رسمی وجود ندارد. نمونه‌ی شناخته‌شده‌ی این مورد، بسته‌های رمزگشایی \en{DVD} (مانند \cmd{libdvdcss}) است که در برخی حوزه‌های قضایی با محدودیت قانونی مواجه‌اند.

از آنجا که این مخازن مستقل از پیکربندی پیش‌فرض توزیع عمل می‌کنند، کاربر باید برای استفاده از آن‌ها، آدرس مخزن را به‌صورت دستی در فایل‌های پیکربندی سیستم مدیریت بسته (\en{Package Manager}) اضافه کند.

\subsection{مدیریت وابستگی‌ها (Dependency Management)}

برنامه‌ها در لینوکس به‌ندرت به‌طور کاملاً مستقل عمل می‌کنند؛ بلکه برای اجرای صحیح خود به سایر اجزای نرم‌افزاری وابسته‌اند. عملیات‌های متداولی مانند ورودی/خروجی (\en{I/O}) معمولاً توسط توابعی انجام می‌شوند که میان چندین برنامه به‌طور مشترک استفاده می‌شوند. این توابع در قالب کتابخانه‌های اشتراکی (\en{Shared Libraries}) نگهداری و در اختیار برنامه‌ها قرار می‌گیرند.

هرگاه بسته‌ای برای اجرای صحیح خود به یک منبع بیرونی—نظیر یکی از این کتابخانه‌ها—نیاز داشته باشد، گفته می‌شود که آن بسته دارای وابستگی (\en{Dependency}) است. تمامی سیستم‌های مدیریت بسته‌ی مدرن دارای سازوکاری برای حل وابستگی (\en{Dependency Resolution}) هستند؛ به این ترتیب که هنگام نصب یک برنامه، پیش‌نیازهای آن نیز به‌طور خودکار شناسایی و نصب می‌شوند.

\subsection{ابزارهای مدیریت بسته: سطح پایین در برابر سطح بالا}

سیستم‌های مدیریت بسته معمولاً از دو لایه‌ی ابزاری تشکیل شده‌اند:

\begin{itemize}
  \item \textbf{ابزارهای سطح پایین (\en{Low-level Tools}):} این ابزارها به‌طور مستقیم با فایل‌های بسته سروکار دارند و وظایفی نظیر نصب، حذف، یا استخراج فیزیکی فایل‌ها را بر عهده می‌گیرند؛ اما به‌تنهایی قادر به حل خودکار وابستگی‌ها نیستند. نمونه‌های شناخته‌شده‌ی این دسته \cmd{dpkg} و \cmd{rpm} هستند.

  \item \textbf{ابزارهای سطح بالا (\en{High-level Tools}):} این ابزارها بر پایه‌ی ابزارهای سطح پایین ساخته شده‌اند و امکاناتی نظیر جست‌وجو در مخازن، پردازش فراداده (\en{Metadata})، و به‌ویژه حل خودکار وابستگی‌ها را فراهم می‌کنند. \cmd{apt} و \cmd{dnf} از نمونه‌های شاخص این دسته‌اند.
\end{itemize}

در این فصل، هم ابزارهای مورد استفاده در توزیع‌های مبتنی بر دبیان (مانند اوبونتو و مشتقات آن) و هم ابزارهای رایج در خانواده‌ی توزیع‌های ردهت را بررسی خواهیم کرد.

در میان توزیع‌های خانواده‌ی ردهت، اگرچه همگی از ابزار سطح پایین یکسانی به نام \cmd{rpm} بهره می‌برند، اما در انتخاب ابزار سطح بالا تنوع بیشتری دیده می‌شود. تمرکز اصلی این نوشتار بر ابزار سطح بالای \cmd{dnf} خواهد بود که امروزه در فدورا (\en{Fedora}) و ردهت اینترپرایز لینوکس (\en{RHEL}) به‌عنوان استاندارد شناخته می‌شود. سایر توزیع‌های این خانواده نیز ابزارهای سطح بالایی با قابلیت‌های مشابه ارائه می‌دهند (مطابق جدول زیر).

\begin{table}[htbp]
\centering
\caption{ابزارهای مدیریت بسته در توزیع‌های شاخص}
\label{tab:pkg-tools}
\begin{tabularx}{\textwidth}{>{\raggedright\arraybackslash}X l l}
\toprule
\textbf{خانواده توزیع} & \textbf{ابزارهای سطح پایین} & \textbf{ابزارهای سطح بالا} \\
\midrule
مبتنی بر دبیان & \cmd{dpkg} & \cmd{apt}، \cmd{apt-get}، \cmd{aptitude} \\
\addlinespace
فدورا، سنت‌اواس استریم (\en{CentOS Stream})، ردهت اینترپرایز لینوکس (\en{RHEL 8} و بالاتر) & \cmd{rpm} & \cmd{dnf}، \cmd{yum} \\
\bottomrule
\end{tabularx}
\end{table}

\begin{note}
\textbf{نکته:} در نسخه‌های قدیمی‌تر خانواده‌ی ردهت (از جمله \en{RHEL 7})، ابزار سطح بالای استاندارد \cmd{yum} بود. از فدورا ۲۲ و \en{RHEL 8} به بعد، \cmd{dnf} به‌عنوان نسل بعدی \cmd{yum} و جانشین رسمی آن معرفی شد و اکنون در تمامی توزیع‌های فعال این خانواده استاندارد محسوب می‌شود.
\end{note}

\section{عملیات رایج مدیریت بسته‌ها}

بسیاری از فرآیندهای مدیریتی را می‌توان از طریق ابزارهای خط فرمان (\en{CLI}) به انجام رساند. در این بخش، به بررسی متداول‌ترینِ این عملیات می‌پردازیم. شایان ذکر است که ابزارهای سطح پایین علاوه بر مدیریت، قابلیت ساخت بسته‌های نرم‌افزاری را نیز دارا هستند که بررسی آن در چارچوب این کتاب نمی‌گنجد.

در مباحث پیش‌رو، عبارت \file{package\_name} به نام سیستمی بسته اشاره دارد که یک شناسهٔ منحصربه‌فرد و داخلی برای مدیر بسته است و نباید با \file{package\_file} که نام فایلیِ فیزیکی بسته است، اشتباه گرفته شود. نام فایل بسته صرفاً نامی است که فایل روی دیسک دارد و معمولاً (اما نه لزوماً) برای سهولت کاربر، شامل نام بسته به همراه اطلاعات نسخه و معماری می‌شود، در حالی که نام سیستمی بسته درون خود فایل تعریف شده و مدیر بسته برای شناسایی و مدیریت آن از این نام استفاده می‌کند. همچنین، پیش از اجرای هرگونه عملیات، پرس‌وجو از مخازن جهت همگام‌سازی پایگاه داده محلی الزامی است. ابزار \cmd{dnf} در خانواده ردهت، این فرآیند را به‌صورت خودکار مدیریت کرده و در صورت انقضای داده‌های محلی، آن‌ها را به‌روزرسانی می‌کند. در مقابل، ابزار \cmd{apt} در خانواده دبیان نیازمند اجرای صریح دستور \cmd{update} برای به‌روزرسانی پایگاه داده محلی است. اگرچه در نمونه‌های ذیل، دستور \cmd{apt update} پیش از هر عملیاتی درج شده، اما در محیط عملیاتی، یکبار اجرای آن در فواصل زمانی چند ساعته کفایت می‌کند.

از آنجا که نصب یا حذف نرم‌افزار یک تغییر در سطح کل سیستم محسوب می‌شود، فارغ از نوع ابزار مورد استفاده، اجرای این فرامین مستلزم برخورداری از امتیازات ابرکاربر (\en{superuser}) است.

\subsection{پایش و جست‌وجوی بسته‌ها در مخازن}

با بهره‌گیری از ابزارهای سطح بالا جهت جستجو در فراداده‌های (\en{Metadata}) موجود در مخازن، امکان مکان‌یابی بسته‌های نرم‌افزاری بر اساس نام یا شرح عملکرد آن‌ها فراهم است (مطابق جدول زیر).

\begin{table}[htbp]
\centering
\caption{فرامین جست‌وجوی بسته در مخازن}
\label{tab:pkg-search}
\begin{tabularx}{\textwidth}{l X}
\toprule
\textbf{خانواده توزیع} & \textbf{دستور (یا فرامین) عملیاتی} \\
\midrule
مبتنی بر دبیان & \cmd{apt update; apt search search\_string} \\
\addlinespace
مبتنی بر ردهت & \cmd{dnf search search\_string} \\
\bottomrule
\end{tabularx}
\end{table}

به‌عنوان مثال، برای جستجوی نرم‌افزار \cmd{emacs} در مخازن از طریق ابزار \cmd{dnf} در سیستم‌های خانواده ردهت، دستور زیر به‌کار می‌رود:

\begin{latin}
\begin{lstlisting}
dnf search emacs
\end{lstlisting}
\end{latin}

\subsection{نصب نرم‌افزار از مخازن رسمی}

ابزارهای مدیریت بسته‌ی سطح بالا، فرآیند دریافت و نصب نرم‌افزار از مخازن را به‌طور کامل خودکارسازی کرده و تمامی وابستگی‌های مورد نیاز (\en{Dependencies}) را پیش از نصب نهایی شناسایی و مستقر می‌کنند (مطابق جدول زیر).

\begin{table}[htbp]
\centering
\caption{فرامین نصب بسته‌های نرم‌افزاری}
\label{tab:pkg-install}
\begin{tabularx}{\textwidth}{l X}
\toprule
\textbf{خانواده توزیع} & \textbf{دستور (یا فرامین) عملیاتی} \\
\midrule
مبتنی بر دبیان & \cmd{apt update; apt install package\_name} \\
\addlinespace
مبتنی بر ردهت & \cmd{dnf install package\_name} \\
\bottomrule
\end{tabularx}
\end{table}

برای نمونه، جهت نصب نرم‌افزار \cmd{emacs} در سیستم‌های مبتنی بر دبیان با استفاده از ابزار \cmd{apt} از دستور زیر استفاده می‌شود:

\begin{latin}
\begin{lstlisting}
apt update; apt install emacs
\end{lstlisting}
\end{latin}

توصیه می‌شود همواره پیش از نصب بسته‌های جدید در سیستم‌های دبیان‌محور، از دستور \cmd{apt update} استفاده کنید تا مدیر بسته فهرست بسته‌ها و آخرین نسخه‌های موجود در مخازن را دریافت و به‌روز کند.

\subsection{نصب مستقیم بسته‌های نرم‌افزاری}

چنانچه فایل بسته‌ای را از منبعی خارج از مخازن رسمی دریافت کرده باشید، امکان نصب مستقیم آن با استفاده از ابزارهای سطح پایین فراهم است. در این شیوه، برخلاف ابزارهای سطح بالا، فرآیند حل خودکار وابستگی‌ها (\en{Dependency Resolution}) انجام نمی‌شود (مطابق جدول زیر).

\begin{table}[htbp]
\centering
\caption{فرامین سطح پایین نصب بسته}
\label{tab:pkg-install-lowlevel}
\begin{tabularx}{\textwidth}{l X}
\toprule
\textbf{خانواده توزیع} & \textbf{دستور عملیاتی} \\
\midrule
مبتنی بر دبیان & \cmd{dpkg -i package\_file} \\
\addlinespace
مبتنی بر ردهت & \cmd{rpm -i package\_file} \\
\bottomrule
\end{tabularx}
\end{table}

به‌عنوان مثال، اگر فایل بسته‌ی \file{emacs-22.1-7.fc7-i386.rpm} از منبعی غیر از مخازن رسمی دانلود شده باشد، با دستور زیر نصب می‌گردد:

\begin{latin}
\begin{lstlisting}
rpm -i emacs-22.1-7.fc7-i386.rpm
\end{lstlisting}
\end{latin}

\begin{note}
\textbf{نکته:} از آنجا که در این روش از ابزار سطح پایین \cmd{rpm} استفاده می‌شود، هیچ‌گونه مکانیزمی برای مدیریت پیش‌نیازها وجود ندارد. در صورتی که سیستم با فقدان وابستگی‌های نرم‌افزاری مواجه شود، فرآیند نصب متوقف شده و ابزار با اعلام خطا خارج می‌گردد.
\end{note}

\subsection{حذف بسته‌های نرم‌افزاری}

فرآیند حذف نرم‌افزار از سیستم می‌تواند هم از طریق ابزارهای سطح بالا و هم به‌وسیله‌ی ابزارهای سطح پایین صورت پذیرد. مزیت استفاده از ابزارهای سطح بالا که در جدول زیر فهرست شده‌اند، مدیریت هوشمندانه‌تر و یکپارچگی بیشتر با کل سیستم است.

\begin{table}[htbp]
\centering
\caption{فرامین حذف بسته}
\label{tab:pkg-remove}
\begin{tabularx}{\textwidth}{l X}
\toprule
\textbf{خانواده توزیع} & \textbf{دستور عملیاتی} \\
\midrule
مبتنی بر دبیان & \cmd{apt remove package\_name} \\
\addlinespace
مبتنی بر ردهت & \cmd{dnf erase package\_name} \\
\bottomrule
\end{tabularx}
\end{table}

برای نمونه، جهت پاکسازی بسته‌ی \cmd{emacs} از یک سیستم مبتنی بر دبیان، از فرمان زیر استفاده می‌شود:

\begin{latin}
\begin{lstlisting}
apt remove emacs
\end{lstlisting}
\end{latin}

در توزیع‌های مبتنی بر دبیان، دستور \cmd{remove} فایل‌های اجرایی را حذف می‌کند اما فایل‌های پیکربندی را باقی می‌گذارد. اگر مایلید تمامی آثار بسته، شامل فایل‌های پیکربندی نیز به‌طور کامل حذف شوند، می‌توانید از پارامتر \cmd{purge} استفاده کنید.

\subsection{ارتقای بسته‌های نرم‌افزاری از مخازن رسمی}

متداول‌ترین وظیفه در فرآیند مدیریت بسته، به‌روزرسانی سیستم جهت بهره‌مندی از آخرین نسخه‌ی بسته‌های نرم‌افزاری است. ابزارهای سطح بالا قادرند این عملیات حیاتی را به‌صورت یکپارچه و در کمترین مراحل به انجام برسانند (جدول زیر).

\begin{table}[htbp]
\centering
\caption{فرامین به‌روزرسانی و ارتقای سیستم}
\label{tab:pkg-upgrade}
\begin{tabularx}{\textwidth}{l X}
\toprule
\textbf{خانواده توزیع} & \textbf{دستور (یا فرامین) عملیاتی} \\
\midrule
مبتنی بر دبیان & \cmd{apt update; apt upgrade} \\
\addlinespace
مبتنی بر ردهت & \cmd{dnf update} \\
\bottomrule
\end{tabularx}
\end{table}

به‌عنوان نمونه، برای اعمال تمامی به‌روزرسانی‌های موجود بر روی بسته‌های نصب‌شده در یک سیستم خانواده دبیان، از ترکیب فرامین زیر استفاده می‌شود:

\begin{latin}
\begin{lstlisting}
apt update; apt upgrade
\end{lstlisting}
\end{latin}

در این فرآیند، ابتدا فهرست آخرین تغییرات از مخازن دریافت شده (\cmd{update}) و سپس نسخه‌های جدید جایگزین نسخه‌های قدیمی می‌شوند (\cmd{upgrade}). در سیستم‌های مبتنی بر ردهت، ابزار \cmd{dnf} هر دو مرحله را به‌صورت خودکار در قالب یک دستور واحد به انجام می‌رساند.

\subsection{ارتقای نرم‌افزار از طریق فایل بسته}

چنانچه نسخه‌ی جدیدی از یک بسته نرم‌افزاری را از منبعی خارج از مخازن رسمی دریافت کرده باشید، می‌توانید با استفاده از ابزارهای سطح پایین، آن را جایگزین نسخه‌ی پیشین مستقر در سیستم نمایید (مطابق جدول زیر).

\begin{table}[htbp]
\centering
\caption{فرامین سطح پایین ارتقای بسته}
\label{tab:pkg-upgrade-lowlevel}
\begin{tabularx}{\textwidth}{l X}
\toprule
\textbf{خانواده توزیع} & \textbf{دستور عملیاتی} \\
\midrule
مبتنی بر دبیان & \cmd{dpkg -i package\_file} \\
\addlinespace
مبتنی بر ردهت & \cmd{rpm -U package\_file} \\
\bottomrule
\end{tabularx}
\end{table}

به‌عنوان مثال، جهت ارتقای نرم‌افزار \cmd{emacs} موجود در یک سیستم ردهت به نسخه‌ی مندرج در فایل \file{emacs-22.1-7.fc7-i386.rpm} از دستور زیر استفاده می‌شود:

\begin{latin}
\begin{lstlisting}
rpm -U emacs-22.1-7.fc7-i386.rpm
\end{lstlisting}
\end{latin}

\begin{note}
\textbf{نکته:} برخلاف ابزار \cmd{rpm} که برای تمایز میان نصب اولیه و ارتقای بسته از گزینه‌های متفاوتی استفاده می‌کند، ابزار \cmd{dpkg} در خانواده دبیان پارامتر اختصاصی برای ارتقاء ندارد و در هر دو حالت از گزینه \cmd{-i} (مخفف \en{install}) استفاده می‌نماید.
\end{note}

\subsection{استخراج فهرست بسته‌های نصب‌شده}

جدول زیر فرامین لازم جهت مشاهده و تهیه لیست کامل تمامی بسته‌های نرم‌افزاری‌ای که هم‌اکنون بر روی سیستم مستقر شده‌اند را ارائه می‌دهد. این دستورها با مراجعه به پایگاه داده محلی مدیر بسته، فهرستی از بسته‌های نصب‌شده را تولید می‌کنند و به مخازن نرم‌افزاری متصل نمی‌شوند.

\begin{table}[htbp]
\centering
\caption{فرامین فهرست‌بندی بسته‌ها}
\label{tab:pkg-list}
\begin{tabularx}{\textwidth}{l X}
\toprule
\textbf{خانواده توزیع} & \textbf{دستور عملیاتی} \\
\midrule
مبتنی بر دبیان & \cmd{dpkg -l} \\
\addlinespace
مبتنی بر ردهت & \cmd{rpm -qa} \\
\bottomrule
\end{tabularx}
\end{table}

در سیستم‌های مبتنی بر دبیان، خروجی دستور \cmd{dpkg -l} افزون بر نام بسته‌ها، وضعیت نصب (\en{Package State}) شامل وضعیت فعلی، وضعیت مطلوب یا انتخاب (\en{Selection State}) و پرچم‌های (\en{Flags}) بسته، و همچنین شرح مختصری (\en{Description}) از هر بسته را نیز نمایش می‌دهد. در مقابل، در خانوادهٔ ردهت، دستور \cmd{rpm -qa} با مراجعه به پایگاه دادهٔ محلی سیستم، فهرستی فشرده شامل نام (\en{Name})، نسخه (\en{Version})، انتشار (\en{Release}) و معماری (\en{Architecture}) تمامی بسته‌هایی را که هم‌اکنون روی سیستم نصب شده‌اند، ارائه می‌کند.

\subsection{استعلام وضعیت نصب بسته}

جدول زیر ابزارهای سطح پایین مورد نیاز جهت اعتبارسنجی وضعیت استقرار یک بسته‌ی مشخص در سیستم را ارائه می‌دهد.

\begin{table}[htbp]
\centering
\caption{فرامین استعلام وضعیت بسته}
\label{tab:pkg-status}
\begin{tabularx}{\textwidth}{l X}
\toprule
\textbf{خانواده توزیع} & \textbf{دستور عملیاتی} \\
\midrule
مبتنی بر دبیان & \cmd{dpkg -s package\_name} \\
\addlinespace
مبتنی بر ردهت & \cmd{rpm -q package\_name} \\
\bottomrule
\end{tabularx}
\end{table}

به‌عنوان مثال، برای اطمینان از نصب بودن بسته‌ی \cmd{emacs} بر روی یک سیستم خانواده دبیان، می‌توان از فرمان زیر استفاده کرد:

\begin{latin}
\begin{lstlisting}
dpkg --status emacs
\end{lstlisting}
\end{latin}

خروجی این دستورات علاوه بر تایید نصب، اطلاعاتی نظیر نسخه‌ی دقیق، وابستگی‌ها و حجم بسته را نیز در اختیار شما قرار می‌دهد. در سیستم‌های ردهت، چنانچه بسته نصب نباشد، ابزار \cmd{rpm} با پیام \en{``package is not installed''} به شما پاسخ خواهد داد.

\subsection{مشاهده جزئیات بسته‌های نصب‌شده}

در صورتی که نام یک بسته‌ی مستقر در سیستم را بدانید، می‌توانید با استفاده از فرامین مندرج در جدول زیر، اطلاعات کاملی از آن—شامل نام، نسخه، نگهدارنده (\en{Maintainer})، وابستگی‌ها، و شرح فنی بسته—را مشاهده کنید.

\begin{table}[htbp]
\centering
\caption{فرامین نمایش اطلاعات تفصیلی بسته}
\label{tab:pkg-info}
\begin{tabularx}{\textwidth}{l X}
\toprule
\textbf{خانواده توزیع} & \textbf{ابزار عملیاتی} \\
\midrule
مبتنی بر دبیان & \cmd{apt show package\_name} \\
\addlinespace
مبتنی بر ردهت & \cmd{dnf info package\_name} \\
\bottomrule
\end{tabularx}
\end{table}

برای نمونه، جهت مطالعه‌ی مستندات و توضیحات مربوط به بسته‌ی \cmd{emacs} در توزیع‌های خانواده دبیان، دستور زیر به‌کار می‌رود:

\begin{latin}
\begin{lstlisting}
apt-cache show emacs
\end{lstlisting}
\end{latin}

\subsection{ردیابی منشأ فایل‌های سیستمی}

برای تشخیص اینکه یک فایل خاص توسط کدام بسته‌ی نرم‌افزاری در سیستم مستقر شده است، می‌توان از دستورات تخصصی مندرج در جدول زیر بهره گرفت.

\begin{table}[htbp]
\centering
\caption{فرامین شناسایی بسته‌ی مرجعِ فایل}
\label{tab:pkg-owner}
\begin{tabularx}{\textwidth}{l X}
\toprule
\textbf{خانواده توزیع} & \textbf{دستور عملیاتی} \\
\midrule
مبتنی بر دبیان & \cmd{dpkg -S file\_name} \\
\addlinespace
مبتنی بر ردهت & \cmd{rpm -qf file\_name} \\
\bottomrule
\end{tabularx}
\end{table}

به‌عنوان مثال، جهت شناسایی بسته‌ای که مسئول نصب فایل اجرایی \file{/usr/bin/vim} در یک سیستم خانواده ردهت است، فرمان زیر اجرا می‌شود:

\begin{latin}
\begin{lstlisting}
rpm -qf /usr/bin/vim
\end{lstlisting}
\end{latin}

این قابلیت به‌ویژه در زمان عیب‌یابی سیستم یا هنگامی که قصد دارید بدانید یک فایل ناشناخته متعلق به کدام ابزار است، بسیار حیاتی و کاربردی خواهد بود.

\section{قالب‌های بسته‌بندی مستقل از توزیع}

در سال‌های اخیر، ارائه‌دهندگان لینوکس قالب‌های بسته‌بندیِ مستقل از توزیع (\en{Universal Packaging Formats}) را معرفی کرده‌اند که به توزیع خاصی وابسته نیستند. این قالب‌ها شامل \cmd{Snap} (توسعه‌یافته توسط کانونیکال)، \cmd{Flatpak} (که ریشه در پروژه‌ی \en{xdg-app} دارد و با مشارکت گسترده‌ی جامعه‌ی گنوم و ردهت شکل گرفت) و \cmd{AppImage} هستند. هرچند سازوکار اجرایی هر یک از این قالب‌ها متفاوت است، اما هدف مشترک آن‌ها تجمیع نرم‌افزار به‌همراه تمامی وابستگی‌هایش در قالب یک بسته‌ی واحد و آماده‌ی نصب است. این ایده چندان نوین نیست؛ در سال‌های آغازینِ لینوکس (دهه‌ی ۱۹۹۰ میلادی)، تکنیکی به نام «اتصال ایستا» (\en{Static Linking}) وجود داشت که فایل اجرایی و تمامی کتابخانه‌های موردنیازش را در یک فایل باینریِ حجیم ترکیب می‌کرد.

این رویکرد مزایای قابل‌توجهی دارد که مهم‌ترین آن‌ها، تسهیل فرآیند توزیع نرم‌افزار است. به‌جای انطباق برنامه با کتابخانه‌ها و فایل‌های سیستمیِ متغیر در هر توزیع، نرم‌افزار تنها یک‌بار ساخته می‌شود و بر روی هر سیستمی قابل نصب خواهد بود. علاوه بر این، برخی از این قالب‌ها—به‌ویژه \cmd{Snap} و \cmd{Flatpak}—جهت ارتقای سطح امنیت، برنامه را در یک محیط ایزوله (\en{Sandbox}) اجرا می‌کنند.

با این حال، معایب جدی نیز بر این رویکرد مترتب است. برنامه‌هایی که با این روش بسته‌بندی می‌شوند، حجم بسیار بالایی دارند که دو چالش اساسی ایجاد می‌کند: نخست، اشغال فضای ذخیره‌سازیِ زیاد، و دوم، افزایش زمان بارگذاری (\en{Load Time}) به‌دلیل حجم بالای فایل. اگرچه این مسائل در سخت‌افزارهای مدرن کمتر محسوس است، اما در سیستم‌های قدیمی‌تر یک معضل واقعی تلقی می‌شود. چالش فنی دیگر، عدم یکپارچگیِ کامل با توزیع میزبان است؛ از آنجا که این برنامه‌ها وابستگی‌های اختصاصی خود را حمل می‌کنند، از زیرساخت‌های مشترک سیستم استفاده نمی‌کنند و گاهی در محیط‌های ایزوله، برای دسترسی به منابع سیستمی جهت عملکرد بهینه با محدودیت مواجه می‌شوند.

در نهایت، ابعاد فلسفیِ این موضوع نیز حائز اهمیت است. به نظر می‌رسد ذی‌نفعان اصلیِ این بسته‌های «همه‌کاره»، توسعه‌دهندگان نرم‌افزارهای انحصاری (\en{Proprietary}) هستند؛ چرا که می‌توانند نسخه‌ی لینوکسیِ محصول خود را تنها یک‌بار تولید کرده و بدون نیاز به بومی‌سازی برای توزیع‌های مختلف، عرضه کنند.

از آنجا که این قالب‌ها پاسخی به نیاز مستقیم کاربران نبوده‌اند و نقش چندانی در تقویت جوامع متن‌باز ایفا نمی‌کنند، تا زمان رفع چالش‌های عملکردی و ساختاری، استفاده‌ی گسترده از آن‌ها توصیه نمی‌شود.

\section{جمع‌بندی}

در فصل‌های آتی، به بررسی طیف گسترده‌ای از ابزارها در حوزه‌های کاربردی گوناگون خواهیم پرداخت. اگرچه بخش عمده‌ای از این برنامه‌ها به‌صورت پیش‌فرض در سیستم موجود هستند، اما گاهی ناگزیر به نصب بسته‌های تکمیلی خواهیم بود. با اتکا به دانش جدیدی که در این فصل کسب کردید و با درک عمیق از سازوکار مدیریت بسته، فرآیند نصب و نگهداری نرم‌افزارهای مورد نیاز دیگر چالشی دشوار نخواهد بود.

\section{ابطال انگاره دشواری نصب نرم‌افزار در لینوکس}

کاربرانی که از سایر پلتفرم‌ها به دنیای لینوکس مهاجرت می‌کنند، گاه با این تصور نادرست مواجه می‌شوند که نصب نرم‌افزار در این سیستم‌عامل فرآیندی پیچیده است و تنوع قالب‌های بسته‌بندی در توزیع‌های مختلف، سدی در برابر کاربر محسوب می‌شود. واقعیت این است که این تنوع تنها برای عرضه‌کنندگان نرم‌افزارهای انحصاری (\en{Proprietary}) که قصد توزیع نسخه‌های باینری و بسته را دارند، یک مانع به شمار می‌آید.

زیست‌بوم نرم‌افزاری لینوکس بر پایه‌ی فلسفه‌ی متن‌باز بنا شده است. زمانی که توسعه‌دهنده‌ای کد منبع برنامه‌ی خود را منتشر می‌کند، نگهدارندگان توزیع‌های مختلف آن را دریافت، کامپایل و برای مخازن اختصاصی خود بسته‌بندی می‌کنند. این رویکرد تضمین می‌کند که نرم‌افزار با ساختار آن توزیع کاملاً سازگار باشد. همچنین، این مدل به کاربر اجازه می‌دهد تا از مزیت «مرکزیت توزیع» بهره‌مند شود و به جای جستجوی پراکنده در وب‌سایت‌های مختلف، تمامی نیازهای خود را از یک منبع واحد و امن تامین کند؛ ایده‌ای که امروزه پلتفرم‌های مالکیتی بزرگ نیز با ایجاد فروشگاه‌های نرم‌افزاری (\en{App Stores}) سعی در تقلید از آن دارند.

مدیریت درایورهای سخت‌افزاری در لینوکس رویکردی متفاوت دارد؛ این درایورها به جای آنکه به‌عنوان بسته‌هایی مجزا در مخازن نرم‌افزاری ارائه شوند، به صورت درونی در هسته لینوکس (\en{Kernel}) ادغام شده‌اند. به بیان دقیق، در اکوسیستم لینوکس مفهومی به نام «دیسک درایور» وجود ندارد. پشتیبانی از یک قطعه سخت‌افزاری، مستقیماً به حضور کد آن در هسته وابسته است. جالب است بدانید هسته لینوکس به‌صورت پیش‌فرض از طیف وسیع‌تری از سخت‌افزارها نسبت به ویندوز پشتیبانی می‌کند، اما اگر سخت‌افزار خاص شما شناسایی نشود، باید علل فنی آن را ریشه‌یابی کرد.

\begin{table}[htbp]
\centering
\caption{علل عدم پشتیبانی درایور در هسته}
\label{tab:driver-issues}
\begin{tabularx}{\textwidth}{c X c}
\toprule
\textbf{علت} & \textbf{تبیین فنی} & \textbf{اولویت} \\
\midrule
تازگی بیش از حد سخت‌افزار & از آنجا که بسیاری از تولیدکنندگان سخت‌افزار، به‌صورت فعال در توسعه‌ی لینوکس مشارکت نمی‌کنند، وظیفه‌ی نگارش درایور بر عهده‌ی جامعه‌ی متن‌باز می‌افتد که این فرآیند مهندسی معکوس و توسعه، زمان‌بر است. & ۱ \\
\addlinespace
تخصصی بودن سخت‌افزار & توزیع‌های مختلف، هسته‌های خود را با پیکربندی‌های متفاوتی کامپایل می‌کنند. از آنجا که هسته لینوکس بسیار انعطاف‌پذیر است، ممکن است یک توزیع خاص برای کاهش حجم هسته، درایور برخی قطعات خاص را حذف کرده باشد. در این صورت، کاربر می‌تواند با دریافت کد منبع درایور، آن را شخصاً کامپایل و نصب نماید. & ۲ \\
\addlinespace
عدم انتشار مستندات فنی (\en{Closed Source}) & برخی فروشندگان نه تنها کد منبع درایور را منتشر نمی‌کنند، بلکه از ارائه‌ی مستندات لازم برای توسعه‌دهندگان نیز خودداری می‌کنند. این پنهان‌کاری مانع از ایجاد درایورهای استاندارد می‌شود. در فرهنگ لینوکس، توصیه می‌شود جهت حفظ پایداری و امنیت، از خرید چنین محصولاتی اجتناب کنید. & ۳ \\
\bottomrule
\end{tabularx}
\end{table}

\section{منابع مطالعه تکمیلی}

توصیه می‌شود زمانی را به شناخت عمیق‌تر سازوکارهای مدیریت بسته در توزیع لینوکس خود اختصاص دهید. هر توزیع، مستندات اختصاصی و جامعی برای ابزارهای مدیریتی خود ارائه کرده است. علاوه بر مستندات رسمی، منابع معتبر زیر برای درک مبانی فنی و مقایسه‌ای بسیار سودمند خواهند بود:

فصل مربوط به مدیریت بسته در پرسش‌های متداول \en{Debian GNU/Linux}، مروری کلی بر این موضوع در سیستم‌های دبیان ارائه می‌دهد:

\begin{latin}
\url{https://www.debian.org/doc/manuals/debian-faq/pkg-basics.en.html}
\end{latin}

صفحه اصلی پروژه \en{RPM}:

\begin{latin}
\url{http://www.rpm.org/}
\end{latin}

مقاله ویکی‌پدیا در مورد فراداده (\en{Metadata}) برای درک مفاهیم پیش‌زمینه‌ای:

\begin{latin}
\url{http://en.wikipedia.org/wiki/Metadata}
\end{latin}

مقاله‌ای مناسب برای مقایسه فرمت‌های \en{Snap}، \en{Flatpak} و \en{AppImage}:

\begin{latin}
\url{https://www.baeldung.com/linux/snaps-flatpak-appimage}
\end{latin}